\NativeTitle{37}{第五单元素养评价（B）}{（多边形的认识）}{90}{6}{一 & 二 & 三 & 四 & 五 & 六}{ &  &  &  &  & }
\NativeSlot{280}{42}{%
\ExamGroup{一、填空。{\normalfont\kaishu\fontsize{10.8bp}{14bp}\selectfont （每空2分，共26分）}}

% p037_s01

\ExamQuestion{1.}{5}%
一个三角形的两个内角分别是\(37^\circ\)和\(49^\circ\)，这是一个\mbox{（\hspace{1.8em}）}三角形；一个三角形的三个内角的度数都相等，它是\mbox{（\hspace{1.8em}）}三角形，也是\mbox{（\hspace{1.8em}）}三角形。\par
}
\NativeSlot{356}{56}{%
% p037_s02

\ExamQuestion{2.}{5}%
窗户打开后，可以用窗户挡风撑\(AB\)将其固定，窗户就不会被风吹得晃动，这里运用到三角形的\mbox{（\hspace{1.8em}）}性。\par
\BookFigure{figures/p037_s02-01.tex}

% p037_s03

\ExamQuestion{3.}{5}%
古人用篱笆围田，用一根长24米的篱笆围成一个等腰梯形，上底是4米，下底是8米，腰是\mbox{（\hspace{1.8em}）}米。\par
}
\NativeSlot{456}{45}{%
% p037_s04

\ExamQuestion{4.}{5}%
右图中，有\mbox{（\hspace{1.8em}）}个平行四边形。\par
\BookFigure{figures/p037_s04-01.tex}

% p037_s05

\ExamQuestion{5.}{5}%
木匠把一块三角形木板分成两个小三角形木块，每个小三角形的内角和是\mbox{（\hspace{1.8em}）}。\par
}
\NativeSlot{532.0}{43}{%
% p037_s06

\ExamQuestion{6.}{5}%
一个等腰梯形的下底是6厘米，上底是3厘米，如果把上底向一个方向延长3厘米，则变成一个\mbox{（\hspace{1.8em}）}形；如果把上底缩短3厘米，则变成一个\mbox{（\hspace{1.8em}）}形。\par
}
\NativeSlot{605.0}{44}{%
% p037_s07

\ExamQuestion{7.}{5}%
手工课上，同学们用木条制作三角形装饰框架。其中两根木条分别长8厘米和11厘米，第三根木条最长是\mbox{（\hspace{1.8em}）}厘米，最短是\mbox{（\hspace{1.8em}）}厘米。（取整厘米数）\par
}
\NativeSlot{680.0}{64}{%
% p037_s08

\ExamQuestion{8.}{5}%
在图内的交叉点中再选一个点，如果要使这四个点能连成一个平行四边形，这个点的选法共有\mbox{（\hspace{1.8em}）}种；如果要连成一个梯形，这个点的选法共有\mbox{（\hspace{1.8em}）}种。\par
\BookFigure{figures/p037_s08-01.tex}
}
